\documentclass[11pt]{amsart}
\usepackage[a4paper,margin=25mm]{geometry}
\usepackage{amssymb,microtype,booktabs}
\usepackage[hidelinks]{hyperref}
\newtheorem{theorem}{Theorem}[section]\newtheorem{proposition}[theorem]{Proposition}\newtheorem{lemma}[theorem]{Lemma}
\newtheorem{corollary}[theorem]{Corollary}\newtheorem{conjecture}[theorem]{Conjecture}
\theoremstyle{definition}\newtheorem{definition}[theorem]{Definition}\theoremstyle{remark}\newtheorem{remark}[theorem]{Remark}
\newcommand{\sinc}{\operatorname{sinc}}\newcommand{\E}{\mathbb E}\newcommand{\cx}{\le_{\mathrm{cx}}}
\newcommand{\todo}[1]{\par\noindent\textbf{[TODO: #1]}\par}
\begin{document}
\title[Two-point extremals with fixed marginals]{Two-point extremals for characteristic-function dependence\\ with fixed marginals}
\author{Alexandr Martinevski}
\address{Independent Researcher, Lithuania}\email{\href{mailto:alex.martinevski@gmail.com}{alex.martinevski@gmail.com}}
\thanks{ORCID: \href{https://orcid.org/0009-0000-4230-4414}{0009-0000-4230-4414}.}
\date{September 2026}
\begin{abstract}
For random variables $X_1,\dots,X_d$ we study the largest characteristic-function defect $|\E e^{i\sum_jt_jX_j}-\prod_j\E e^{it_jX_j}|$. In every regime where we determine the exact constant,
the extremal couplings make the phase sum two-point, and a certificate made of two tangent lines proves optimality. Without restrictions on the marginals, a companion paper establishes the sharp benchmark
$C_d=\max_{0\le x\le\pi}\{\cos^d(x/d)-\cos x\}$ ($C_3=2/\sqrt3$, $C_4=9/7$) from an exact profile in the moduli of the means. Here the marginal laws are prescribed. For
a common symmetric unimodal law $F$ we prove $K_d^F=\max_x\{\varphi(x/(d\,\E|X|))^d-\cos x\}$ under explicit conditions --- for the uniform and triangular laws for
all $d\ge3$, for the B-spline and Epanechnikov laws for $d\ge4$ --- and show that, under strict versions of these conditions, attainment is equivalent to complete mixability of $|X|$;
Gaussian marginals do not attain the bound. The uniform law is the least restrictive symmetric unimodal marginal for $d\ge4$ but not for $d=3$; without symmetry the law $\frac14\delta_0+\frac34U[0,1]$ beats it
for every $d\ge4$ and is asymptotically extremal among all unimodal laws, through the sharp inequality $\mathrm{Var}X\ge\frac97(\E|X-\mathrm{med}X|)^2$. For every non-degenerate
law with bounded support that is lattice or satisfies Cram\'er's condition, $K_d^F=2-\pi^2\mathrm{Var}X/(2d\,(\E|X-\mathrm{med}X|)^2)+O(d^{-3/2})$: the median replaces the centre, and no symmetry,
unimodality or mixability is needed. Without symmetry the best two-point split is asymmetric; that it is optimal is conjectured. The finite verifications are short tables of explicit inequalities, checked by monotone brackets, and by interval arithmetic for the non-symmetric comparison.
\end{abstract}
\maketitle

\section{Introduction}\label{sec:intro}
For random variables $X_1,\dots,X_d$ and frequencies $t\in\mathbb R^d$ consider the \emph{characteristic-function defect}
\[\Delta=\E e^{i\sum_jt_jX_j}-\prod_j\E e^{it_jX_j}.\]
It vanishes under independence, and its vanishing at all $t$ on the diagonal is sub-independence; it is the basic quantity behind characteristic-function
measures of dependence. We ask how large $|\Delta|$ can be, in three settings of decreasing freedom: (a) free laws with $|Z_j|\le1$ in place of $e^{it_jX_j}$;
(b) prescribed moduli $|\E Z_j|=r_j$; (c) prescribed marginal laws $X_j\sim F$ (or $F_j$), with the coupling and the frequencies free.

\medskip\noindent\textbf{Results.} (a)--(b) are settled in the companion paper \cite{I}: the sharp profile $C_d(\vec r)=\prod r_j-\cos(\min\{\sum\arccos r_j,\pi\})$ and the constant
$C_d=\max_{0\le x\le\pi}\{\cos^d(x/d)-\cos x\}$ (recalled as Theorem~\ref{thm:profile} and Corollary~\ref{cor:Cd}). (c): for symmetric unimodal $F$,
\[K_d^F=\kappa_d^F:=\max_x\{\varphi(x/(d\,\E|X|))^d-\cos x\}\]
under explicit conditions (Theorem~\ref{thm:main}), for all $d\ge3$ for the uniform and triangular laws and for $d\ge4$ for the B-spline and Epanechnikov laws
(Table~\ref{tab:const}); the analogue for different laws holds under joint mixability (Theorem~\ref{thm:het}). Under strict versions of the hypotheses, attainment is equivalent to
complete mixability of $|X|$ (Theorem~\ref{thm:att}); Gaussian marginals do not attain $\kappa_d^G$ (Corollary~\ref{cor:gauss}). The uniform law has the largest constant
among symmetric unimodal marginals for $d\ge4$ but not for $d=3$ (Theorems~\ref{thm:ext} and~\ref{thm:d3}); without symmetry it is beaten for every $d\ge4$ by
$\frac14\delta_0+\frac34U[0,1]$, which is asymptotically extremal among all unimodal laws (Theorem~\ref{thm:finite-d}, Corollary~\ref{cor:asym-ext}), by the sharp inequality
$\mathrm{Var}X\ge\frac97(\E|X-\mathrm{med}X|)^2$ for unimodal laws, a moment form of the known worst-case expected shortfall (Proposition~\ref{prop:97}). For every non-degenerate law with bounded support that is lattice or satisfies Cram\'er's condition,
$K_d^F=2-\pi^2\mathrm{Var}X/(2d\,(\E|X-\mathrm{med}X|)^2)+O(d^{-3/2})$ (Theorem~\ref{thm:price}); the factor $\mathrm{Var}X/(\E|X-\mathrm{med}X|)^2\ge1$ is the price of
fixing the marginals. Without symmetry, the best two-point split is asymmetric; its optimality is Conjecture~\ref{conj:nonsym}.

\medskip\noindent\textbf{The two-point principle.} In every exact regime proved here the extremal phase sum takes two values, and optimality is certified by two
tangent lines; for non-symmetric laws the same structure gives the lower bound (Theorem~\ref{thm:nonsym-lb}) and is conjectured to be optimal:
\begin{center}\footnotesize\begin{tabular}{p{3.2cm}p{4.3cm}p{4.6cm}}\toprule
setting & the two points & certificate\\\midrule
fixed moduli & $\pm\sum\arccos r_j$ (common sign) & concavity of $\log\cos$\\
symmetric laws & $\pm t\sum\E|X_j|$ (sign, mixing) & a tent at $\pm x_0$ (Lemma~\ref{lem:window})\\
general laws (lower bound) & means of the two quantile parts & two tangents (Lemma~\ref{lem:twotan})\\
scales, unimodal class & two scales & majorant of $\sinc$ (Lemma~\ref{lem:maj})\\
large $d$ & odd multiples of $\pi$ & odd-lattice threshold (Lemma~\ref{lem:thr})\\
unimodal moment problem & two atoms of Khinchine's $Z$ & two-point decomposition (Lemma~\ref{lem:dec})\\\bottomrule
\end{tabular}\end{center}
The mechanism is the same throughout: a single linear constraint survives, extreme points are two-point, and equality in the certificate forces the
coupling to be a mix of the two parts --- which is where complete and joint mixability (Wang--Wang, Puccetti--Wang--Wang) enter.

\medskip\noindent\textbf{Organization.} Section~\ref{sec:prelim} collects tools; Section~\ref{sec:free} recalls the free-law results of \cite{I}; Section~\ref{sec:twopoint} develops the two-point principle;
Sections~\ref{sec:main}--\ref{sec:att} treat fixed marginals; Section~\ref{sec:nonsym} records the non-symmetric phenomenon; Section~\ref{sec:unimodal} identifies the extremal unimodal marginal without symmetry;
Section~\ref{sec:open} lists open problems. The finite
verifications are in Appendix~\ref{app:fin}.

\section{Preliminaries}\label{sec:prelim}
\begin{definition}
For laws $F_1,\dots,F_d$ with characteristic functions $\varphi_j$,
\[K(F_1,\dots,F_d)=\sup\Bigl|\E e^{i\sum_jt_jX_j}-\prod_j\varphi_j(t_j)\Bigr|\]
over $t\in\mathbb R^d$ and couplings with $X_j\sim F_j$; $K_d^F=K(F,\dots,F)$. For symmetric $F$ with $0<m=\E|X|<\infty$, $\kappa_d^F=\sup_{x>0}\{\varphi(x/(dm))^d-\cos x\}$.
\end{definition}
\noindent For symmetric $F$ we write $s_0\in(0,\infty]$ for the first positive zero of $\varphi$ ($s_0=\infty$ if there is none; first lobe $[0,s_0)$), $s^*=\sup_{s\ge s_0}|\varphi(s)|$, $\gamma=f(0)m$ when $F$ has a density $f$, and
$x_c=2.3311\ldots$ for the root of $\cos x+x\sin x=1$ in $(\frac\pi2,\pi)$.

\medskip\noindent\textbf{Convex order and mixability.} $Y\cx V$ means $\E g(Y)\le\E g(V)$ for all convex $g$. A $d$-tuple $(F_1,\dots,F_d)$ is \emph{jointly mixable} (JM) if some coupling
has $\sum_jX_j$ constant; $F$ is $d$-\emph{completely mixable} ($d$-CM) if $(F,\dots,F)$ is JM. We use three criteria:
\begin{enumerate}
\item[(M1)] (Wang--Wang 2011 \cite{WW11}) A law with a nonincreasing density on $[0,b]$ is $d$-CM iff its mean is at least $b/d$.
\item[(M2)] (Wang--Wang 2016 \cite{WW16}, Thms~2.1(a) and~3.2) Laws with nonincreasing densities on $[0,l_j]$ and means $\mu_j$ are JM iff $\sum_j\mu_j\ge\max_jl_j$.
\item[(M3)] (Puccetti--Wang--Wang 2012 \cite{PWW12}, Thm~4.2) A law with a concave density on a bounded interval is $d$-CM for every $d\ge3$.
\end{enumerate}
\noindent\textbf{Unimodality.} By Khinchine's theorem (see \cite{DJ88}) a symmetric unimodal $X$ is $CW$ with $W\sim U[-1,1]$ independent of $C\ge0$, so
$\varphi(s)=\E\sinc(sC)$; if $X$ has a density, $|X|$ has a nonincreasing density on $[0,b]$, $b=\operatorname{ess\,sup}C$, and $m=\E C/2\le b/2$.
\begin{lemma}[the first lobe]\label{lem:lobe}
$\sinc z\ge1-z/\pi$ for all $z\ge0$. Hence for symmetric unimodal $F$, $\varphi(s)\ge1-2sm/\pi$; for $d\ge3$ the interval $[0,\pi/(dm)]$ lies in the first lobe,
$\varphi(\pi/(dm))\ge1-2/d>0$, and $\kappa_d^F\ge1+P_d(\pi)>1$, where $P_d(x)=\varphi(x/(dm))^d$.
\end{lemma}
\begin{proof}On $[0,\pi]$ let $q(z)=\sin z-z(\pi-z)/\pi$. Then $q(\pi-z)=q(z)$, $q(0)=q'(0)=0$ and $q'(\frac\pi2)=0$, and $q''=\frac2\pi-\sin z$ is
positive on $[0,z_1)$ and negative on $(z_1,\frac\pi2]$, $z_1=\arcsin\frac2\pi$. On $[0,z_1]$, $q$ is convex with a double zero at $0$, so $q\ge0$; on $[z_1,\frac\pi2]$, $q$ is
concave, so $q\ge\min\{q(z_1),q(\frac\pi2)\}\ge0$ ($q(\frac\pi2)=1-\frac\pi4$). By symmetry $q\ge0$ on $[0,\pi]$. On $[\pi,2\pi]$, $\sin z\ge-(z-\pi)\ge-z(z-\pi)/\pi$. For $z\ge2\pi$, $1-z/\pi\le-1\le\sinc z$. Then
$\varphi(s)=\E\sinc(sC)\ge1-s\E C/\pi$, and $\E C=2m$.\end{proof}
\begin{lemma}[lattice separation]\label{lem:latsep}
For all real $y$: $1-\cos y\ge\frac2{\pi^2}\,\mathrm{dist}(y,2\pi\mathbb Z)^2$ and $1+\cos y\ge\frac2{\pi^2}\,\mathrm{dist}(y,\pi(2\mathbb Z+1))^2$.
\end{lemma}
\begin{proof}With $\delta$ the distance, $|\delta|\le\pi$ and $1-\cos\delta=2\sin^2(\delta/2)\ge2(\delta/\pi)^2$, since $\sin u\ge2u/\pi$ on $[0,\frac\pi2]$; and $1+\cos y=1-\cos(y-\pi)$.\end{proof}

\section{Free laws: the profile and the constant $C_d$}\label{sec:free}
For complex random variables with $|Z_j|\le1$ let $\Delta=\E\prod_jZ_j-\prod_j\E Z_j$. Put $r_j=|\E Z_j|$, $\theta_j=\arccos r_j\in[0,\frac\pi2]$, $R=\prod_jr_j$, $s=\sum_j\theta_j$, and
\[C_d(r_1,\dots,r_d)=\sup\{|\Delta|:\ |Z_j|\le1,\ |\E Z_j|=r_j\}.\]
The following results are proved in the companion paper \cite{I}; we recall them because the regimes with fixed marginals below are measured against them.
\begin{theorem}[profile {\cite[Thm~2.1]{I}}]\label{thm:profile}
For every $d\ge1$ and $r\in[0,1]^d$: $\ C_d(r_1,\dots,r_d)=R-\cos(\min\{s,\pi\})$. For $s\le\pi$ it is attained by $Z_j=e^{i\varepsilon\theta_j}$ with a common fair sign
$\varepsilon$; for $s\ge\pi$ (value $1+R$) by a law with $\prod_jZ_j\equiv-1$.
\end{theorem}
\begin{corollary}[the constant $C_d$ {\cite[Cor.~2.2, Thm~2.4]{I}}]\label{cor:Cd}
The supremum of $|\E\prod Z_j-\prod\E Z_j|$ over all $|Z_j|\le1$ is
\[C_d=\max_{0\le x\le\pi}\{\cos^d(x/d)-\cos x\};\qquad C_2=1,\quad C_3=2/\sqrt3,\quad C_4=9/7,\]
and for $d\ge3$ it is attained only when all $\arccos|\E Z_j|$ equal $x^*_d/d$, where $x^*_d\in(0,\pi)$ is the unique maximizer.
\end{corollary}
\begin{proof}Only the last assertion needs a comment: by Theorem~\ref{thm:profile} and strict concavity of $\log\cos$ on $[0,\frac\pi2)$, $\prod\cos\theta_j\le\cos^d(s/d)$ with
equality iff the $\theta_j$ are equal, and the maximizer $x^*_d=d\theta^*_d$ is unique and interior by \cite[Thm~2.4]{I}.\end{proof}

\section{The two-point principle}\label{sec:twopoint}
Throughout this section $X_1,\dots,X_d\sim F$, $t\in\mathbb R^d$, $S=\sum_jt_jX_j$ and $z=\prod_j\varphi(t_j)$. The extremal couplings in every problem solved exactly in this paper
make the phase sum two-point; this section collects the three facts behind that: a lower bound by two-point fusions, a dual certificate made of two
tangent lines, and the fact that equality forces complete mixability.

\subsection{Relaxation}
Let $U$ be uniform on $(0,1)$ and $g_j$ the quantile function of $\mathrm{sgn}(t_j)X$. The \emph{comonotone sum} is $V=\sum_j|t_j|\,g_j(U)=h(U)$, with $h$
nondecreasing. For symmetric $F$, $g_j=F^{-1}$ and $V=AX$ with $A=\sum_j|t_j|$.
\begin{lemma}[comonotone domination]\label{lem:como}
For every coupling, $S\cx V$. Consequently $|\E e^{iS}-z|\le R(V,z):=\sup\{|\E e^{iT}-z|:\ T\cx V\}$.
\end{lemma}
\begin{proof}Standard (Meilijson--N\'adas \cite{MN79}; Dhaene et al.\ \cite{Dh02}): $\E(S-k)_+\le\E(V-k)_+$ for all $k$ and $\E S=\E V$.\end{proof}

\subsection{Quantile-split fusions and the lower bound}
\begin{definition}\label{def:fusion}
For $p\in(0,1)$ let $a_p=\E[V\mid U<p]$, $b_p=\E[V\mid U\ge p]$, $v_p=h(p)$. The \emph{fusion} $T_p$ equals $a_p$ on $\{U<p\}$ and $b_p$ on $\{U\ge p\}$; it satisfies
$T_p\cx V$, and its characteristic value is $w_p=p\,e^{ia_p}+(1-p)\,e^{ib_p}$.
\end{definition}
For a law $F$ and $p\in(0,1)$ let $F_-^{(p)}$ and $F_+^{(p)}$ be the laws of $X_-^{(p)}=F^{-1}(pU)$ and $X_+^{(p)}=F^{-1}(p+(1-p)U)$ (the lower and upper $p$-parts).
\begin{theorem}[fusions are attained under mixability of the parts]\label{thm:fusion-lb}
Let $t_1=\dots=t_d=t>0$. If $F_-^{(p)}$ and $F_+^{(p)}$ are $d$-completely mixable, there is a coupling with $S=T_p$ (where $V=dtX$); hence
$K_d^F\ge|w_p-\varphi(t)^d|$.
\end{theorem}
\begin{proof}Take an event $E$ with $\mathbb P(E)=p$, independent of two complete mixes $(L_1,\dots,L_d)$ of $F_-^{(p)}$ and $(R_1,\dots,R_d)$ of $F_+^{(p)}$, and put
$X_j=L_j$ on $E$ and $X_j=R_j$ off $E$. Then $X_j\sim pF_-^{(p)}+(1-p)F_+^{(p)}=F$, and $t\sum_jX_j$ equals $dt\,\E L_1=a_p$ on $E$ and $b_p$ off $E$.\end{proof}
\begin{corollary}[symmetric laws]\label{cor:sym-lb}
Let $F$ be symmetric and $|X|$ be $d$-completely mixable. Then $K_d^F\ge\varphi(t)^d-\cos(dt\,\E|X|)$ for every $t$; hence $K_d^F\ge\kappa_d^F$. By
Wang--Wang, this holds if $|X|$ has a nonincreasing density on $[0,b]$ and $d\ge b/\E|X|$.
\end{corollary}
\begin{proof}For $p=\frac12$ the parts are the laws of $-|X|$ and $|X|$, so $a_{1/2}=-dt\,\E|X|$, $b_{1/2}=dt\,\E|X|$ and $w_{1/2}=\cos(dt\,\E|X|)$. (Equivalently:
a common random sign times a complete mix of $|X|$.)\end{proof}
\begin{remark}If $F$ has a concave density on a bounded interval, both parts have concave densities for every $p$, hence are $d$-CM for $d\ge3$
(Puccetti--Wang--Wang, Thm~4.2): every fusion is attained. This is the entry point of Sections~\ref{sec:nonsym} and~\ref{sec:unimodal}.
\end{remark}

\subsection{Two tangents}
\begin{lemma}[two-tangent certificate]\label{lem:twotan}
Let $f\in C^1(\mathbb R)$, $p\in(0,1)$, and let $\ell_a,\ell_b$ be the tangent lines to $f$ at $a=a_p$ and $b=b_p$, with $\ell_a'\le\ell_b'$ and $\ell_a(v_p)=\ell_b(v_p)$. If
$G=\max(\ell_a,\ell_b)\ge f$ on the convex hull of the range of $V$, then $\sup_{T\cx V}\E f(T)=p\,f(a)+(1-p)\,f(b)$, attained by $T_p$.
\end{lemma}
\begin{proof}$T\cx V$ takes values in that convex hull and $G$ is convex, so $\E f(T)\le\E G(T)\le\E G(V)$. The lines cross at $v_p$ and $\ell_a'\le\ell_b'$, so $G(V)=\ell_a(V)$ on $\{U<p\}$ (where $V\le v_p$) and
$G(V)=\ell_b(V)$ on $\{U\ge p\}$. By linearity $\E G(V)=p\,\ell_a(a)+(1-p)\,\ell_b(b)=p\,f(a)+(1-p)\,f(b)=\E f(T_p)$.\end{proof}
\begin{lemma}[stationarity]\label{lem:stat}
If $h$ is continuous at $p$, then $\frac{d}{dp}\bigl[p\,f(a_p)+(1-p)f(b_p)\bigr]=\ell_a(v_p)-\ell_b(v_p)$. In particular, at an interior maximizer in $p$ the
tangent lines cross at the split point $v_p$, as required in Lemma~\ref{lem:twotan}; the slope condition $\ell_a'\le\ell_b'$ and the majorization must still be checked.
\end{lemma}
\begin{proof}$\frac{da_p}{dp}=\frac{v_p-a_p}p$ and $\frac{db_p}{dp}=\frac{b_p-v_p}{1-p}$, so the derivative is $f(a)-f(b)+f'(a)(v_p-a)+f'(b)(b-v_p)=\ell_a(v_p)-\ell_b(v_p)$.\end{proof}

\subsection{The symmetric case: the window}
For symmetric $F$ and equal frequencies the fusion at $p=\frac12$ is $\pm x_0$ with $x_0=dt\,\E|X|$, and in the direction of $-1$ Lemma~\ref{lem:twotan} becomes
the classical ``tent''. Let $x_c\approx2.3311$ be the root in $(\frac\pi2,\pi)$ of $\cos x+x\sin x=1$.
\begin{lemma}[window]\label{lem:window}
For $x_0\in[x_c,\pi]$ and all real $s$: $\cos s\ge\cos x_0-\sin x_0\,(|s|-x_0)$, with equality only at $|s|=x_0$ if $x_c<x_0<\pi$.
\end{lemma}
\begin{proof}By symmetry let $s\ge0$ and put $L(s)=\cos x_0-\sin x_0(s-x_0)$, $g=\cos-L$; then $g(x_0)=g'(x_0)=0$ and $g''=-\cos$. On $[\frac\pi2,\frac{3\pi}2]$, $g$ is convex
and tangent to $0$ at $x_0$, so $g\ge0$. On $[0,\frac\pi2]$, $g$ is concave, so its minimum is at an end: $g(\frac\pi2)\ge0$, and $g(0)=1-\cos x_0-x_0\sin x_0\ge0$
because $x\mapsto\cos x+x\sin x$ decreases on $[\frac\pi2,\pi]$ and equals $1$ at $x_c$. For $s\ge\frac{3\pi}2$, $L(s)\le L(\frac{3\pi}2)\le-1$, since
$1+\cos x_0\le\sin x_0(\frac{3\pi}2-x_0)$ is $\cot\frac{x_0}2\le\frac{3\pi}2-x_0$, true on $[x_c,\pi]$ ($\cot\frac{x_c}2<0.43<\frac\pi2$). Strictness away from $x_0$ follows
from strict convexity or concavity on each piece when $x_c<x_0<\pi$.\end{proof}
\begin{corollary}\label{cor:window}
If $F$ is symmetric, $x_0=A\,\E|X|\in[x_c,\pi]$ and $T\cx AX$, then $\E\cos T\ge\cos x_0$, attained by the fusion $\pm x_0$.
\end{corollary}
\begin{proof}$\E\cos T\ge\cos x_0-\sin x_0(\E|T|-x_0)\ge\cos x_0$, since $\E|T|\le\E|AX|=x_0$.\end{proof}

\subsection{Equality forces complete mixability}
\begin{proposition}\label{prop:eq-mix}
Let $t_1=\dots=t_d=t>0$, $V=dtX$, and suppose a coupling gives $S=T_p$ almost surely, i.e.\ $S$ equals $a_p$ on an event $E$ with $\mathbb P(E)=p$ and $b_p$ on
$E^c$. If $F$ is continuous at $F^{-1}(p)$, then on $E$ every $X_j$ has law $F_-^{(p)}$ and on $E^c$ law $F_+^{(p)}$; in particular both parts are $d$-completely
mixable.
\end{proposition}
\begin{proof}(Bathtub principle; $t>0$.) For any event $B$ with $\mathbb P(B)=p$, $\E[X;B]\ge\E[X;X<F^{-1}(p)]=p\,\E X_-^{(p)}$, with equality only if $B=\{X<F^{-1}(p)\}$ up to
null sets. Summing over $j$: $\E[S;E]=t\sum_j\E[X_j;E]\ge p\,dt\,\E X_-^{(p)}=p\,a_p$, and $\E[S;E]=p\,a_p$ by assumption; so equality holds for every $j$.\end{proof}
\begin{corollary}[symmetric case]\label{cor:eq-mix-sym}
If $F$ is symmetric, $t>0$, $x_0=dt\,\E|X|\in(x_c,\pi)$ and a coupling gives $\E\cos S=\cos x_0$ with $S=t\sum X_j$, then $\sum_j|X_j|\equiv d\,\E|X|$ almost surely.
\end{corollary}
\begin{proof}By Lemma~\ref{lem:window} and Corollary~\ref{cor:window}, equality forces $\E|S|=x_0$ and $|S|=x_0$ a.s. Since $|S|\le t\sum|X_j|$ and both sides
have mean $x_0$, $t\sum_j|X_j|\equiv x_0$.\end{proof}

\section{Symmetric unimodal marginals: the main theorem}\label{sec:main}
Let $F$ be symmetric and unimodal with density $f$, $m=\E|X|$, $\gamma=f(0)\,m$, and let $s_0\in(0,\infty]$ be the first positive zero of $\varphi$ (the \emph{first lobe}
is $[0,s_0)$), $s^*=\sup_{s\ge s_0}|\varphi(s)|$ ($s^*=0$ if $s_0=\infty$), and $P_d(x)=\varphi(x/(dm))^d$. For $0<x_0\le\pi$ and $\eta\in[0,\pi]$ put $\hat x_0=\max\{x_0,x_c\}$ and
\[B_d(x_0,\eta)=\min\Bigl\{1+P_d(x_0)\cos\eta,\ \ -\cos\hat x_0+\sin\hat x_0\bigl(x_0-\hat x_0+\min\{\tfrac{\gamma\eta^2}{x_0},\eta\}\bigr)+P_d(x_0)\cos\eta\Bigr\}.\]
\noindent\textbf{Conditions.} (W$_d$) $\ \sup_{0<x_0\le\pi,\ 0\le\eta\le\pi}B_d(x_0,\eta)\le\kappa_d^F$; \quad (T$_d$) $\ 1+s^*\le\kappa_d^F$.
\begin{theorem}[main theorem]\label{thm:main}
Let $F$ be symmetric and unimodal with density, with $\log\varphi$ concave on its first lobe.
\begin{enumerate}
\item[(i)] If (W$_d$) and (T$_d$) hold, then $K_d^F\le\kappa_d^F$.
\item[(ii)] If $F$ is supported on $[-b,b]$ and $d\ge b/m$, then $K_d^F\ge\kappa_d^F$.
\item[(iii)] If $F$ has bounded support and $\gamma<\gamma_0:=\min_{x_c\le x\le\pi}\frac{2x}{\pi^2\sin x}=0.6519\ldots$, then there is $d_F$ with $K_d^F=\kappa_d^F$ for all $d\ge d_F$.
\end{enumerate}
\end{theorem}
\begin{remark}In the window $x_0\ge x_c$, if $\gamma\sin x_0\,\eta^2/x_0\le P_d(x_0)(1-\cos\eta)$ for all $\eta$, the second term of $B_d$ is at most $P_d(x_0)-\cos x_0\le\kappa_d^F$: this is
the two-tangent certificate of Section~\ref{sec:twopoint} in every direction (``exactness''). The finite verifications of the Appendix check (W$_d$) by
accepting each box either by a margin or by exactness.\end{remark}

\begin{lemma}[shift of a unimodal law]\label{lem:shift}
If $V$ is symmetric and unimodal with density $g$, then $\E|V-\eta|\le\E|V|+\min\{g(0)\eta^2,|\eta|\}$ for all $\eta$.
\end{lemma}
\begin{proof}$\E|V-\eta|-\E|V|=\int_0^{|\eta|}(2F_V(u)-1)\,du$, and $0\le2F_V(u)-1\le\min\{2g(0)u,1\}$ for $u\ge0$ because $g\le g(0)$.\end{proof}
\begin{lemma}[clamped tent]\label{lem:clamped}
Let $V$ be symmetric and unimodal with density $g$, $x_0=\E|V|\in(0,\pi]$, $\gamma_V=g(0)\,x_0$, and $S'\cx V$. Then for every real $\eta$,
\[\E\cos(S'-\eta)\ \ge\ \cos\hat x_0-\sin\hat x_0\bigl(x_0-\hat x_0+\min\{\gamma_V\eta^2/x_0,|\eta|\}\bigr).\]
For $V=TX$ with $T>0$, $\gamma_V=\gamma$.
\end{lemma}
\begin{proof}Lemma~\ref{lem:window} at $\hat x_0\in[x_c,\pi]$, applied pointwise at $s=S'-\eta$, gives $\cos(S'-\eta)\ge\cos\hat x_0-\sin\hat x_0(|S'-\eta|-\hat x_0)$. Take expectations,
use $\E|S'-\eta|\le\E|V-\eta|$ (convex order) and Lemma~\ref{lem:shift}. For $V=TX$, $g(0)=f(0)/T$ and $x_0=Tm$.\end{proof}

\begin{proof}[Proof of Theorem~\ref{thm:main}]\emph{(i) Reductions.} Replacing $X_j$ by $-X_j$ does not change the law, so let $t_j\ge0$; put $T=\sum_jt_j$, $x_0=Tm$ and
$P=\prod_j\varphi(t_j)$ (real). By Lemma~\ref{lem:como}, $|\E e^{iS}-P|\le R(T,P):=\sup\{|w-P|:\ w\in W_T\}$, $W_T=\{\E e^{iS'}:S'\cx TX\}$; the map $P\mapsto R(T,P)$ is convex,
$R(T,0)\le1$, and $W_T$ is invariant under conjugation.

\emph{Tail.} If some $t_j\ge s_0$, then $|P|\le s^*$ and the defect is at most $1+s^*\le\kappa_d^F$ by (T$_d$).

\emph{First lobe.} Otherwise all $t_j<s_0$ and $0<P\le P_d(x_0)$ by concavity of $\log\varphi$. By convexity, $R(T,P)\le\max\{R(T,0),R(T,P_d(x_0))\}\le\max\{1,R(T,P_d(x_0))\}$,
and $1<\kappa_d^F$ (Lemma~\ref{lem:lobe}); so it suffices to bound $R(T,P_d(x_0))$.

\emph{$x_0>\pi$.} Since $\log\varphi$ is concave with $(\log\varphi)'(0)=0$, $\varphi$ decreases on the first lobe, which contains $[\pi/(dm),x_0/(dm)]$ (Lemma~\ref{lem:lobe}); so $R\le1+P_d(x_0)\le1+P_d(\pi)$, the value of the formula at $x=\pi$, hence $\le\kappa_d^F$.

\emph{$0<x_0\le\pi$.} For $w\in W_T$, $|w-P|=\sup_\eta\{-\mathrm{Re}(e^{-i\eta}w)+P\cos\eta\}$ (direction $\pi+\eta$), and by conjugation invariance $\eta\in[0,\pi]$ suffices. Now
$-\mathrm{Re}(e^{-i\eta}w)=-\E\cos(S'-\eta)$ is at most $1$ and, by Lemma~\ref{lem:clamped}, at most $-\cos\hat x_0+\sin\hat x_0(x_0-\hat x_0+\min\{\gamma\eta^2/x_0,\eta\})$. Hence
$R(T,P_d(x_0))\le\sup_\eta B_d(x_0,\eta)\le\kappa_d^F$ by (W$_d$).

\emph{(ii)} $|X|$ has a nonincreasing density on $[0,b]$ with mean $m\ge b/d$, so it is $d$-completely mixable (Wang--Wang); apply Corollary~\ref{cor:sym-lb}.

\emph{(iii)} Let $b<\infty$ bound the support. For $d\ge b/m$, (ii) applies. As $d\to\infty$, $P_d\to1$ uniformly on $[0,\pi]$ and $\kappa_d^F\ge1+P_d(\pi)\to2$, so (T$_d$) holds
since $s^*<1$. In (W$_d$) with $x_0\ge x_c$: since $(1-\cos\eta)/\eta^2\ge2/\pi^2$ on $(0,\pi]$ and $\gamma<\gamma_0$, we get $\gamma\sin x_0\,\eta^2/x_0\le(1-\delta)(1-\cos\eta)$ for some
$\delta>0$, hence exactness once $P_d\ge1-\delta$. For $x_0<x_c$ (where $\hat x_0=x_c$ and the second term is $-1+\sin x_c(x_0+\min\{\gamma\eta^2/x_0,\eta\})+P_d\cos\eta$, since $\cos x_c+x_c\sin x_c=1$):
if $\eta\ge\frac12$, the first term is at most $1+\cos\frac12<1.88$; if $\eta\le\frac12$, then $x_0+\min\{\gamma\eta^2/x_0,\eta\}\le\max\{(1+\gamma)\eta,\ x_c+\gamma\eta^2/x_c\}\le2.41$,
so the second term is at most $-1+2.41\sin x_c+1<1.75$. Hence $\sup B_d<1.88$ there, while $\kappa_d^F\to2$.\end{proof}

\begin{lemma}[one frequency beyond the lobe]\label{lem:onebeyond}
Let $F$ be symmetric, $t_1\ge0$ arbitrary and $t_2,\dots,t_d$ in the first lobe, $R'=\prod_{j\ge2}\varphi(t_j)$, $P=\varphi(t_1)R'$ and $\beta=m\sum_{j\ge2}t_j\le\pi$. Then for
every coupling $|\E e^{iS}-P|\le|\varphi(t_1)|-|P|+2\sin(\beta/2)$.
\end{lemma}
\begin{proof}Let $Z_1=e^{it_1X_1}$, $Y=\sum_{j\ge2}t_jX_j$, $W=e^{iY}$. Then $\E e^{iS}-P=\varphi(t_1)(1-R')+\E[Z_1(W-1)]$, and $0<R'\le1$. Since $|W-1|=2|\sin(Y/2)|\le2g(|Y|)$
with the concave nondecreasing $g(y)=\sin(y/2)$ for $y\le\pi$, $g=1$ beyond, Jensen gives $\E|W-1|\le2g(\E|Y|)\le2\sin(\beta/2)$, as $\E|Y|\le\beta$.\end{proof}
\noindent When (T$_d$) fails, this lemma replaces it: in the region $|P|>\kappa_d^F-1$, log-concavity bounds $R'\le\varphi(\beta/((d-1)m))^{d-1}$, which caps $\beta$, while the lemma
settles small $\beta$ (Theorem~\ref{thm:d3} and Lemma~\ref{lem:c3}).

\begin{theorem}[different laws]\label{thm:het}
Let $F_1,\dots,F_d$ be symmetric unimodal with densities $f_j$, $0<m_j=\E|X_j|<\infty$, $b_j=\operatorname{ess\,sup}|X_j|$, $\log\varphi_j$ concave on the first lobes, $\gamma=\max_jf_j(0)m_j$,
$s^*=\max_js_j^*$, and
\[K^*=\sup_{t\ge0}\Bigl\{\prod_j\varphi_j(t_j)-\cos\sum_jt_jm_j\Bigr\}.\]
(i) If $1+s^*\le K^*$ and $B(x_0,\eta;P(t))\le K^*$ for all $t$ in the first lobes with $x_0=\sum t_jm_j\le\pi$ and all $\eta\in[0,\pi]$ (with $B$ as $B_d$ with $P_d(x_0)$ replaced by
$P(t)=\prod\varphi_j(t_j)$), then $K(F_1,\dots,F_d)\le K^*$. (ii) If $K^*$ is attained at $t^*$ with $\sum_jt_j^*m_j\ge\max_jt_j^*b_j$, then $K(F_1,\dots,F_d)\ge K^*$.
\end{theorem}
\begin{proof}(i) By symmetry of each $F_j$ we may take $t\ge0$. As in Theorem~\ref{thm:main}, now with $V=\sum_jt_jF_j^{-1}(U)$. The quantile function of $V$ is a sum of quantile functions that are convex on $(\frac12,1)$ (for a symmetric unimodal density
$q'(u)=1/f(q(u))$ increases there), so $V$ is symmetric and unimodal, with density $g(0)=1/\sum_jt_j/f_j(0)$ at $0$ and $g(0)\E|V|=\sum t_jm_j/\sum(t_j/f_j(0))\le\gamma$ (mediant inequality); so Lemma~\ref{lem:clamped} holds with $\gamma_V\le\gamma$, and $B$ increases in $\gamma$.
No reduction of unequal frequencies is needed, because $K^*$ is a supremum over all $t$. If $x_0>\pi$, scaling $t$ to $x_0=\pi$ increases every $\varphi_j(t_j)$ on the first lobes, so
$1+P(t)\le K^*$. (ii) The moduli $t_j^*|X_j|$ have nonincreasing densities on $[0,t_j^*b_j]$ with means $t_j^*m_j$; by (M2) they are JM, and a common random sign gives $S=\pm x_0^*$,
hence the defect $P(t^*)-\cos x_0^*=K^*$.\end{proof}

\section{Families and a table of constants}
\begin{corollary}\label{cor:families}
$K_d^F=\kappa_d^F$ for: the uniform law and the triangular law, all $d\ge3$; the cubic B-spline (sum of three $U[-\frac12,\frac12]$) and the Epanechnikov
law, all $d\ge4$ (and $K_3^F\le\kappa_3^F$ for the B-spline). For the uniform law on $[-1,1]$, $\kappa_d=\max_x\{\sinc(2x/d)^d-\cos x\}$;
$\kappa_3=1.0849633414$, $\kappa_4=1.1907222653$.
\end{corollary}
\begin{proof}Theorem~\ref{thm:main}(i)--(ii), with the hypotheses verified in Appendix~\ref{app:fin} (log-concavity, tail, (W$_d$)); where (T$_d$) fails
(uniform, $d=3,4$) the tail is closed by Lemma~\ref{lem:onebeyond} and Lemma~\ref{lem:c3} for one frequency beyond the lobe, and by $|P|\le s^{*2}=0.0472<\kappa_d-1$ for two or more. Mixability: $b/m=2$ (uniform), $3$ (triangular), $48/13$ (B-spline),
$8/3$ (Epanechnikov).\end{proof}
\begin{table}[h]\centering\small
\begin{tabular}{lccl}\toprule
law & $d=3$ & $d=4$ & status\\\midrule
free laws, $C_d$ & $1.1547005384=2/\sqrt3$ & $1.2857142857=9/7$ & sharp\\
uniform $[-1,1]$ & $1.0849633414$ & $1.1907222653$ & $K_d=\kappa_d$, all $d\ge3$\\
triangular & $1.0765878231$ & $1.1642014101$ & $K_d=\kappa_d$, all $d\ge3$\\
cubic B-spline & $1.0808671327$ & $1.1656390998$ & $K_3\le\kappa_3$; $K_d=\kappa_d$, $d\ge4$\\
Epanechnikov & $1.0800186407$ & $1.1758685660$ & $K_d=\kappa_d$, $d\ge4$ ($d=3$ open)\\
$\frac34U[-5,5]+\frac14U[-9,9]$ & $1.0874848676$ & & $K_3=\kappa_3^{F}>\kappa_3$\\
Gaussian & $1.0848101338$ & $1.1630296791$ & $K_d<\kappa_d^G$ (bound not attained)\\\bottomrule
\end{tabular}
\caption{Values of $\kappa_d^F=\max_x\{\varphi(x/(d\,\E|X|))^d-\cos x\}$ (and $C_d$), computed to 30 digits.}\label{tab:const}
\end{table}

\section{The price of fixing the marginals}\label{sec:price}
Let $X$ have support in $[a,b]$, $\mu=\E X$, $m$ a median, $\sigma^2=\mathrm{Var}X>0$, $\tau=\E|X-m|$.
\begin{theorem}\label{thm:price}
If $F$ is lattice or satisfies Cram\'er's condition ($\sup_{|t|\ge T}|\varphi(t)|<1$ for every $T>0$), then
\[K_d^F=2-\frac{\pi^2\sigma^2}{2d\,\tau^2}+O(d^{-3/2})\qquad(d\to\infty).\]
\end{theorem}
\noindent No symmetry, unimodality or mixability is assumed: the median replaces the centre of symmetry. Since $\sigma\ge\E|X-\mu|\ge\tau$, the factor
$\sigma^2/\tau^2\ge1$ is the \emph{price of fixing the marginals}, with equality iff $X$ is two-point with probabilities $\frac12$ (the rate of $C_d$).
\begin{corollary}For symmetric $F$ (bounded support, lattice or Cram\'er), $2-K_d^F=\frac{\pi^2}{2d}\frac{\mathrm{Var}X}{(\E|X|)^2}+O(d^{-3/2})$; for the uniform law
the factor is $4/3$, for the triangular law $3/2$.\end{corollary}
\begin{lemma}[odd lattice]\label{lem:odd}
For integrable $Y$ and $|c|\le\pi$: $\E|Y-c|\ge\pi-|\E Y|-2\,\E\,\mathrm{dist}(Y,\pi(2\mathbb Z+1))$.
\end{lemma}
\begin{proof}Let $Y'$ be a nearest odd multiple of $\pi$, $D=|Y-Y'|$, $\nu=\E Y'$; then $|\nu|\le|\E Y|+\E D$ and $\E|Y-c|\ge\E|Y'-c|-\E D$. From
$|y-c|\ge|y|-c\,\mathrm{sgn}\,y$ and, if $P_+=\mathbb P(Y'>0)\ge\frac12$, $\E|Y'|=2\E Y'_+-\nu\ge2\pi P_+-\nu$, we get $\E|Y'-c|\ge P_+(2\pi-2c)+c-\nu\ge\pi-\nu$;
the case $P_-\ge\frac12$ is symmetric.\end{proof}
\begin{lemma}[threshold]\label{lem:thr}
Let $V=\sum_js_jg_j(U)$ with $s_j\ge0$, $A=\sum s_j$, each $g_j$ the quantile function of $X$ or $-X$. If $\E Y=0$ and $Y\cx V-\E V$, then
$\E\cos Y\ge-1+\frac{\tau^2}{2\pi^2}(\pi/\tau-A)_+^2$.
\end{lemma}
\begin{proof}All $g_j(U)-g_j(\frac12)$ have the sign of $U-\frac12$, so $\mathrm{med}\,V=\sum s_jg_j(\frac12)$ and $\E|V-\mathrm{med}\,V|=A\tau$. With $c=\mathrm{med}\,V-\E V$ we
have $|c|\le A\tau$ and, by convexity, $\E|Y-c|\le A\tau$. If $A\tau<\pi$, Lemma~\ref{lem:odd} gives $\E D\ge(\pi-A\tau)/2$ for $D=\mathrm{dist}(Y,\pi(2\mathbb Z+1))$, and
Lemma~\ref{lem:latsep} gives $\E(1+\cos Y)\ge\frac2{\pi^2}(\E D)^2$.\end{proof}
\begin{proof}[Proof of Theorem~\ref{thm:price}]\emph{Lattice reduction.} If $X\in a+h\mathbb Z$ with maximal span $h$, replacing $t_j$ by $t_j+2\pi/h$ multiplies $\E e^{iS}$ and $z=\prod\varphi(t_j)$
by the same unimodular constant; so we may take $|t_j|\le\pi/h$. In both cases there are $T_0>0$ and $q<1$ with $|\varphi(t)|\le q$ for $T_0\le|t|$ (in the reduced
range), and near $0$: $\log|\varphi(t)|\le-\frac{\sigma^2}2t^2+C_1t^4$, $|\varphi(t)|\le e^{-\sigma^2t^2/4}$, $|\arg\varphi(t)-\mu t|\le C_2|t|^3$ ($|\varphi|^2$ is even).

\emph{Upper bound.} Let $\Delta=|\E e^{iS}-z|$. (i) If some $|t_j|\ge T_0$, $\Delta\le1+q$. (ii) If all $|t_j|<T_0$ and $\sum t_j^2\ge M/d$ with $M>4\pi^2/\tau^2$,
$\Delta\le1+e^{-\sigma^2M/(4d)}\le2-\sigma^2M/(8d)$. (iii) Otherwise $A=\sum|t_j|\le\sqrt M$; write $\E e^{iS}=e^{i\E S}u$ with $u=\E e^{iY}$, $Y=S-\E S\cx V-\E V$
(Lemma~\ref{lem:como}), and $z=\rho\,e^{i(\E S+\varepsilon)}$, $|\varepsilon|\le C_2\sum|t_j|^3=O(d^{-3/2})$. By Lemma~\ref{lem:thr}, $\mathrm{Re}\,u\ge-1+\eta$ with
$\eta=\kappa(\pi/\tau-A)_+^2$, $\kappa=\tau^2/(2\pi^2)$, so $|u-\rho|^2\le(1+\rho)^2-2\rho\eta$ and $\Delta\le1+\rho-\frac{\rho\eta}{1+\rho}+\rho|\varepsilon|$. Also
$\rho\le\exp(-\sigma^2A^2/(2d)+O(d^{-2}))$ since $\sum t_j^2\ge A^2/d$. With $A^*=\pi/\tau$, $x=(A^*-A)_+$ and $\rho\ge\frac12$ (otherwise $\Delta\le\frac32$),
\[\Delta\le1+e^{-\sigma^2A^{*2}/(2d)}\Bigl(1+\tfrac{2\sigma^2A^*x}d\Bigr)-\tfrac\kappa3x^2+O(d^{-3/2})\le2-\frac{\sigma^2A^{*2}}{2d}+O(d^{-3/2}),\]
since $\max_{x\ge0}\{\frac{2\sigma^2A^*}dx-\frac\kappa3x^2\}=\frac{3\sigma^4A^{*2}}{\kappa d^2}=O(d^{-2})$ and $e^{-y}\le1-y+y^2/2$.

\emph{Lower bound} (a fusion with approximate mixing, cf.\ Section~\ref{sec:twopoint}). Split at the median: $F_\mp$ are the laws of $F^{-1}(U/2)$ and $F^{-1}((1+U)/2)$,
with means $\mu_\mp$; then $\mu_+-\mu_-=2\tau$ and $\mu-\mu_-=\tau$. On an event of probability $\frac12$, independent of the uniform variables $U,U'$, put $X_j=F_-^{-1}(\{U+j/d\})$, otherwise
$X_j=F_+^{-1}(\{U'+j/d\})$: each $X_j\sim F$, and since a monotone function's Riemann sums at shifted grids are within its range of the integral, the two
sums are $d\mu_\mp+e_\mp$ with $|e_\mp|\le b-a$. At $t=\pi/(d\tau)$, $e^{itd\mu_+}=e^{itd\mu_-}$ and $\E e^{iS}=e^{itd\mu_-}v$ with $\mathrm{Re}\,v\ge1-t^2(b-a)^2/2$, while
$z=-\rho\,e^{itd\mu_-}e^{i\beta}$ with $\beta=O(d^{-2})$ and $\rho=\exp(-\pi^2\sigma^2/(2d\tau^2)+O(d^{-3}))$. Hence $\Delta\ge\mathrm{Re}(v+\rho e^{i\beta})\ge1+\rho-O(d^{-2})$.
\end{proof}

\section{Extremality of the uniform law}\label{sec:ext}
By Khinchine's theorem a symmetric unimodal $X$ is $CW$ with $W\sim U[-1,1]$ independent of $C\ge0$, so $\varphi(t)=\E\sinc(tC)$, $m=\E C/2$ and $b=\operatorname{ess\,sup}C$; the
uniform law with the same $m$ is $U[-2m,2m]$. With $Y=C/\E C$: $\ \varphi(x/(dm))=\E\sinc(zY)$, $z=2x/d$, and $b\le dm$ means $Y\le d/2$.
\begin{lemma}[majorant of $\sinc$]\label{lem:maj}
Let $\hat h=\sinc$ on $[0,\frac\pi2]$ and $\hat h(z)=\frac4{\pi^2}(\pi-z)$ on $[\frac\pi2,\pi]$. Then $\hat h$ is concave on $[0,\pi]$ and $\hat h\ge\sinc$.
\end{lemma}
\begin{proof}The line is the tangent to $\sinc$ at $\frac\pi2$ ($\sinc\frac\pi2=\frac2\pi$, $\sinc'\frac\pi2=-\frac4{\pi^2}$) and $\sinc$ is concave on $[0,\frac\pi2]$, so $\hat h$ is $C^1$ and
concave. On $[\frac\pi2,\pi]$, $\hat h\ge\sinc$ is $p(z)=\frac4{\pi^2}z(\pi-z)-\sin z\ge0$. Now $p$ is symmetric about $\frac\pi2$ with $p(0)=p(\frac\pi2)=0=p'(\frac\pi2)$, and
$p''=-\frac8{\pi^2}+\sin z$ increases on $[0,\frac\pi2]$: $p$ is concave then convex there, convex near $\frac\pi2$ with a double zero, hence $p\ge0$.\end{proof}
\begin{lemma}[dominance]\label{lem:dom}
If $Y\ge0$, $\E Y=1$, $0<z\le\frac\pi2$ and $zY\le\pi$, then $\E\sinc(zY)\le\sinc z$, with equality only if $Y\equiv1$.
\end{lemma}
\begin{proof}$\E\sinc(zY)\le\E\hat h(zY)\le\hat h(z)=\sinc z$ (Jensen). Equality in Jensen needs $zY$ inside an interval where $\hat h$ is affine and containing $z$; the only
affine piece is $[\frac\pi2,\pi]$, which forces $z=\frac\pi2$ and $Y\ge1$, hence $Y\equiv1$.\end{proof}
\begin{theorem}\label{thm:ext}
Let $d\ge4$ and $F$ symmetric unimodal with $b\le dm$, $\varphi$ nonincreasing on its first lobe and $s^*\le\sinc(2\pi/d)$. Then $\kappa_d^F\le\kappa_d$, with equality
iff $F$ is uniform. Hence, whenever $K_d^F=\kappa_d^F$ (Theorem~\ref{thm:main}), the uniform law has the largest constant among these symmetric unimodal marginals.
\end{theorem}
\begin{proof}For $0<x\le\pi$: $z=2x/d\le\pi/2$ because $d\ge4$, and $zY\le x\le\pi$ because $Y\le d/2$; so $0\le\varphi(x/(dm))\le\sinc(2x/d)$ by Lemma~\ref{lem:dom} ($\sinc\ge0$ on $[0,\pi]$),
and $\varphi(x/(dm))^d-\cos x\le\sinc(2x/d)^d-\cos x\le\kappa_d$. For $x>\pi$: $|\varphi(x/(dm))|\le\max\{\varphi(\pi/(dm)),s^*\}\le\sinc(2\pi/d)$, so the value is at most
$1+\sinc(2\pi/d)^d$, the uniform formula at $\pi$, which is $<\kappa_d$ since its maximizer lies in $(0,\pi)$. Equality forces a maximizer in $(0,\pi)$ and
$Y\equiv1$ there.\end{proof}
\begin{remark}[why $d\ge4$]The majorant $\hat h$ touches $\sinc$ exactly up to $\frac\pi2$, and $2\pi/d\le\frac\pi2$ iff $d\ge4$. For $d=3$ the relevant argument $2x^*/3\approx1.98$
lies beyond $\frac\pi2$, where mixing two scales helps.\end{remark}
\begin{remark}[symmetry is needed]\label{rem:symneeded}
The symmetry assumption in Theorem~\ref{thm:ext} cannot be dropped. The law $F=\frac14\delta_0+\frac34U[0,1]$ is unimodal, and
$\varphi(t)=\frac14+\frac34\frac{e^{it}-1}{it}$ satisfies $|\varphi(t)|<1$ for $t\ne0$ and $|\varphi(t)|\to\frac14$, so Cram\'er's condition holds. Its median is $\frac13$,
$\E|X-\frac13|=\frac7{24}$ and $\mathrm{Var}X=\frac7{64}$, so $\mathrm{Var}X/(\E|X-\mathrm{med}X|)^2=\frac97<\frac43$. By Theorem~\ref{thm:price}, $K_d^F>K_d^U$ for all sufficiently
large $d$. (In fact $\frac97$ is the smallest value of this ratio over unimodal laws, and $K_d^F>K_d^U$ already for every $d\ge4$; see Section~\ref{sec:unimodal}.)
\end{remark}
\begin{theorem}[different laws]\label{thm:ext-het}
Let $d\ge4$ and let $F_1,\dots,F_d$ satisfy Theorem~\ref{thm:het}, with the maximum attained at $t^*$ where joint mixability holds. Then $K(F_1,\dots,F_d)\le\kappa_d$,
with equality iff all $F_j$ are uniform.
\end{theorem}
\begin{proof}Scaling $t$ down to $x_0=\pi$ does not decrease the formula on the first lobes, so $x_0^*=\sum t_j^*m_j\le\pi$. With $u_j=t_j^*m_j$ and $Y_j=C_j/\E C_j$, joint mixability gives
$2u_jY_j=t_j^*C_j\le t_j^*b_j\le x_0^*\le\pi$ and $2u_j\le t_j^*b_j$ (as $b_j\ge2m_j$). Hence $\varphi_j(t_j^*)\le\hat h(2u_j)$ (Jensen), and since $\hat h$ is concave and positive
on $[0,\pi)$, $\log\hat h$ is concave: $\prod_j\hat h(2u_j)\le\hat h(2x_0^*/d)^d=\sinc(2x_0^*/d)^d$. So $K\le\sinc(2x_0^*/d)^d-\cos x_0^*\le\kappa_d$.

\emph{Equality.} $\log\hat h$ is strictly concave on $[0,\pi)$ ($\log\sinc$ on $[0,\frac\pi2]$, $\log(\pi-z)$ on $[\frac\pi2,\pi)$), so $K=\kappa_d$ forces
$u_j=x_0^*/d$ for all $j$. Then $2u_j=2x_0^*/d\le\frac\pi2$ and $\varphi_j(t_j^*)=\E\sinc(2u_jY_j)=\sinc(2u_j)$, so Lemma~\ref{lem:dom} gives $Y_j\equiv1$: each $F_j$ is uniform.
Conversely, for uniform $F_j$ (of any scales) the formula reduces to $\max_x\{\sinc(2x/d)^d-\cos x\}=\kappa_d$ at equal $u_j$, where the moduli are identical uniforms and
hence jointly mixable.\end{proof}
\begin{theorem}[$d=3$: the threshold is sharp]\label{thm:d3}
For $F^*=\frac34U[-5,5]+\frac14U[-9,9]$: $K_3^{F^*}=\kappa_3^{F^*}=1.087484867576\ldots>\kappa_3=1.084963341374\ldots$
\end{theorem}
\begin{proof}\emph{Lower bound.} $|X|$ has a nonincreasing density on $[0,9]$ and mean $3=9/3$, so it is $3$-CM (Wang--Wang); Corollary~\ref{cor:sym-lb} at $x=2.97$ gives the value.
\emph{Upper bound.} Here $m=3$, $\gamma=4/15$, $\log\varphi$ is concave on the first lobe $[0,0.59949)$ (Appendix), and (W$_3$) holds (Appendix: 54 boxes). (T$_3$) fails
($s^*\le0.13537$), so we use Lemma~\ref{lem:onebeyond}: with $\tau=\kappa_3^{F^*}-1$, the region $|P|>\tau$ needs $\varphi(\beta/6)^2\ge R'>\tau/s^*$, i.e.\ $\beta<\beta_{\max}=6\varphi^{-1}(\sqrt{\tau/s^*})=1.0891$,
while $|\varphi(t_1)|-|P|+2\sin(\beta/2)\le s^*-\tau+2\sin(\beta/2)\le\kappa_3^{F^*}$ for $\beta\le\beta_1=1.0932$. Two frequencies beyond the lobe give $|P|\le s^{*2}<\tau$.\end{proof}
\begin{proposition}[optimal $\kappa$-profile over the $d=3$ unimodal class]\label{prop:d3sup}
Over symmetric unimodal laws with $3$-completely mixable $|X|$, $\sup_F\kappa_3^F=\sup_x\{\hat h_x(1)^3-\cos x\}=1.0874848690\ldots$, where $\hat h_x$ is the least concave majorant
of $y\mapsto\sinc(2xy/3)$ on $[0,\frac32]$; it is attained by the normalized scale $Y=C/\E C\in\{0.83346,\frac32\}$ with probabilities $0.75014,0.24986$, and $F^*$ is within $1.4\cdot10^{-9}$.
\end{proposition}
\noindent This concerns the formula $\kappa_3^F$. Whether $\sup_FK_3^F$ over the same class equals this value is open (Section~\ref{sec:open}); it holds for $F^*$
by Theorem~\ref{thm:d3}.
\begin{proof}For fixed $x$ the constraint $\E Y=1$ on $Y\in[0,\frac32]$ is one linear constraint, so $\sup\E\sinc(2xY/3)$ is attained by two-point laws and equals the concave majorant at
$1$ --- the two-point principle once more.\end{proof}

\section{Attainment, Gaussian marginals and the mixability conjecture}\label{sec:att}
\noindent\textbf{Strict conditions.} (W$_d^{\rm s}$): $B_d(x_0,\eta)<\kappa_d^F$ for all $(x_0,\eta)\ne(x_0^*,0)$, where $x_0^*\in(x_c,\pi)$ is the unique maximizer of $P_d(x)-\cos x$;
(T$_d^{\rm s}$): $1+s^*<\kappa_d^F$.
\begin{theorem}[attainment forces mixability]\label{thm:att}
Let $F$ be symmetric unimodal with density, $\log\varphi$ strictly concave on its first lobe, and (W$_d^{\rm s}$), (T$_d^{\rm s}$). If $K_d^F=\kappa_d^F$, then some coupling has
$\sum_j|X_j|\equiv d\,m$. Consequently, under these conditions, $K_d^F=\kappa_d^F$ if and only if $|X|$ is $d$-completely mixable.
\end{theorem}
\begin{proof}Take frequencies $t^{(n)}\ge0$ and couplings with defect $\to\kappa_d^F$. By (T$_d^{\rm s}$) all $t_j^{(n)}<s_0$ eventually. In the proof of Theorem~\ref{thm:main} the defect is at
most $1+P_d(\pi)$ if $x_0>\pi$ (strictly below $\kappa_d^F$ since $x_0^*<\pi$) and, by convexity in $P$, at most $\lambda+(1-\lambda)\sup_\eta B_d(x_0,\eta)$ otherwise, where
$\lambda=1-P/P_d(x_0)\in[0,1)$. Since the defect tends to $\kappa_d^F>1$, (W$_d^{\rm s}$) and compactness give $\lambda^{(n)}\to0$, $x_0^{(n)}\to x_0^*$ and maximizing directions $\eta\to0$;
$\lambda^{(n)}\to0$ and strict concavity of $\log\varphi$ force the frequencies to the equal point $t^*$. Couplings with fixed
marginals form a tight family; a limit coupling at $t^*$ has defect $\kappa_d^F$, attained only in the direction $\eta=0$: $P_d(x_0^*)-\E\cos S=\kappa_d^F$, i.e.\ $\E\cos S=\cos x_0^*$.
Corollary~\ref{cor:eq-mix-sym} gives $\sum_j|X_j|\equiv dm$. Conversely, if $|X|$ is $d$-CM, Corollary~\ref{cor:sym-lb} and Theorem~\ref{thm:main}(i) give $K_d^F=\kappa_d^F$.\end{proof}
\begin{corollary}[Gaussian marginals]\label{cor:gauss}
For standard Gaussian marginals and every $d\ge3$,
\[K_d^G<\kappa_d^G=\max_x\{e^{-\pi x^2/(4d)}-\cos x\};\]
$\kappa_d^G$ is an upper bound that is not attained.
\end{corollary}
\begin{proof}$\varphi=e^{-t^2/2}$ has no zeros ($s^*=0$) and $\log\varphi$ is strictly concave; $\gamma=1/\pi$; (W$_d^{\rm s}$) holds for all $d\ge3$ (Appendix: brackets for $d\le19$; beyond, Lemma~\ref{lem:binf} and
$\kappa_{20}^G=1.6926>1-\cos x_c$ with Lemma~\ref{lem:mono-d}; strict exactness at $x_0^*$, e.g.\ $\gamma\sin x_0^*/(x_0^*P_3(x_0^*))=0.168<4/\pi^2$). Since $|X|$ is unbounded, $\sum|X_j|$ cannot be constant.\end{proof}
\begin{proposition}[the mixability conjecture fails for $d=3$]\label{prop:threepoint}
Let $\mathbb P(X=0)=\mathbb P(X=1)=\mathbb P(X=-1)=\frac13$. Then $|X|$ is $3$-CM (one zero at a uniformly random place) and $\kappa_3^F=(5\sqrt5-7)/4=1.0450849719$, but $K_3^F\ge1.0673306671$.
\end{proposition}
\begin{proof}With probability $\frac13$ let $X_1=X_2=X_3=0$; otherwise let $\varepsilon=\pm1$ be a fair sign and $X_3=\varepsilon$, $X_1=X_2=-\varepsilon$; each $X_j$ has law $F$. At
$t=(\frac{2\pi}5,\frac{2\pi}5,\frac{5\pi}6)$, $S=\varepsilon(t_3-t_1-t_2)=\varepsilon\,\frac\pi{30}$ off the event $\{X_1=X_2=X_3=0\}$, so $\E e^{iS}=\frac13+\frac23\cos\frac\pi{30}$, while $\prod\varphi(t_j)=\varphi(\frac{2\pi}5)^2\varphi(\frac{5\pi}6)<0$ with
$\varphi(t)=\frac13+\frac23\cos t$; the difference is $1.0673306671$.\end{proof}
\noindent Outside the unimodal class, attainment of the formula is therefore not characterized by mixability; a compensating coupling with frequencies of
different sizes beats it. For $d\ge4$ the question is open (Section~\ref{sec:open}).

\section{A new phenomenon without symmetry}\label{sec:nonsym}
For a law without symmetry, the fusions of Section~\ref{sec:twopoint} at a split $p\ne\frac12$ put an imaginary part into the fused point, and this can increase the
defect. For equal frequencies $t>0$ (by conjugation, $t<0$ gives nothing new) let $D(t,p)=|p\,e^{i\,dt\mu_-(p)}+(1-p)e^{i\,dt\mu_+(p)}-\varphi(t)^d|$, where $\mu_\mp(p)$ are the means of
the lower and upper $p$-parts.
\begin{theorem}\label{thm:nonsym-lb}
If $F$ has a concave density on a bounded interval, then $K_d^F\ge\sup_{t,p}D(t,p)$ for every $d\ge3$.
\end{theorem}
\begin{proof}Both parts have concave densities, hence are $d$-completely mixable for $d\ge3$ (Puccetti--Wang--Wang, Thm~4.2); apply Theorem~\ref{thm:fusion-lb}.\end{proof}
\noindent For Beta$(1,2)$ (density $2(1-x)$ on $[0,1]$) and $d=3$: $\sup D=1.0726392817$ at $t^*=5.18958$, $p^*=0.43813$, while the best symmetric split
$p=\frac12$ gives $1.0681$. The optimal two-point split is asymmetric.
\begin{conjecture}\label{conj:nonsym}If $F$ has a concave density on a bounded interval (possibly with analogues of the conditions of Theorem~\ref{thm:main}), then $K_d^F=\sup_{t,p}D(t,p)$ for $d\ge3$.\end{conjecture}
\noindent Evidence: at $t^*$ the comonotone relaxation lies within $4\cdot10^{-7}$ of $D(t^*)$ (explicit two-tangent certificates in $3600$ directions) for
Beta$(1,2)$ ($d=3,4$) and an asymmetric triangular law ($d=3$). The upper-bound tools (an equal-variable reduction of frequency splits, a monotone-path
principle) will appear in a sequel.

\section{Without symmetry: the extremal unimodal marginal}\label{sec:unimodal}
By Theorem~\ref{thm:price}, for large $d$ the least restrictive marginal minimizes $\mathrm{Var}X/\tau^2$, $\tau=\E|X-\mathrm{med}X|$; among symmetric unimodal
laws the minimum is $\frac43$, attained by the uniform law. Without symmetry it is smaller (Remark~\ref{rem:symneeded}), and the uniform law is not extremal at any $d\ge3$.

\noindent\textbf{Unimodality.} $X$ is unimodal about $M$ if its distribution function is convex on $(-\infty,M)$ and concave on $(M,\infty)$ (an atom at $M$ is allowed);
$M$ is then a mode. By Khinchine's theorem \cite{Kh38,DJ88} this holds iff $X=M+UZ$ with $U\sim U[0,1]$ independent of $Z$.
\begin{proposition}[the $9/7$ inequality]\label{prop:97}
Let $X$ be unimodal with $0<\mathrm{Var}X<\infty$ and $\tau=\E|X-\mathrm{med}X|$.
\begin{enumerate}
\item[(i)] $\mathrm{Var}X\ge\frac97\,\tau^2$, with equality if and only if $X$ is an affine image of $\frac14\delta_0+\frac34U[0,1]$.
\item[(ii)] If some median of $X$ is a mode of $X$ (in particular, if $X$ is symmetric), then $\mathrm{Var}X\ge\frac43\,\tau^2$, with equality if and only if $X$ is uniform.
\end{enumerate}
\end{proposition}
\noindent With $\mathrm{ES}_{1/2}(X)=2\int_{1/2}^1F^{-1}(u)\,du$ one has $\tau=\mathrm{ES}_{1/2}(X)-\E X$, so (i) is the case $p=\frac12$ of the worst-case expected shortfall
over unimodal laws with given mean and variance, which is known in closed form: Li, Shao, Wang and Yang \cite{LSWY18} treat Range Value-at-Risk under mean, variance
and unimodality; see also \cite{BKV20,ZBYS26}, and \cite[Remark~8]{ZY25} for the formula $\sup(\mathrm{ES}_p-\E X)=\sigma\sqrt{p(8-9p)}/(3(1-p))$, $p\le\frac12$, which equals
$\frac{\sqrt7}3\sigma$ at $p=\frac12$. We include a short proof by the two-point principle; it also gives the uniqueness of the extremal law and part (ii), which we have not
found stated. Part (ii) does not follow from the Gauss--Winckler inequality $\E(X-M)^2\ge\frac43(\E|X-M|)^2$ about a mode $M$ \cite{Wi1866}, since
$\mathrm{Var}X=\E(X-M)^2-(\E X-M)^2$.
\begin{lemma}[two-point decomposition]\label{lem:dec}
Let $\nu$ be a probability law on $\mathbb R$ and $g$ measurable with values in $[0,1]$, $\int g\,d\nu=c$. Then $\nu$ is a mixture of laws with at most two atoms, each with
$\int g=c$.
\end{lemma}
\begin{proof}Let $\mu$ be the law of $g(Z)$, $Z\sim\nu$, with restrictions $\mu_\pm$ to $\{y\gtrless c\}$ and $D=\int(y-c)\,d\mu_+=\int(c-y')\,d\mu_-$. Then
\[\mu=\mu(\{c\})\delta_c+\iint\frac{(c-y')\delta_y+(y-c)\delta_{y'}}{y-y'}\cdot\frac{(y-y')\,\mu_+(dy)\,\mu_-(dy')}D,\]
a mixture of two-point laws with mean $c$ (the marginals are $\mu_+$ and $\mu_-$). Disintegrating $\nu$ along $g(Z)$ and drawing $Z$ independently from the two
conditional laws lifts each two-point law of $g(Z)$ to a mixture of two-point laws of $Z$ with the same $g$-mean.\end{proof}
\begin{proof}[Proof of Proposition~\ref{prop:97}]\emph{Reduction.} Fix a mode $M$ of $X$ and write $X=M+UZ$. The ratio $\mathrm{Var}X/\tau^2$ is invariant under
$X\mapsto\alpha X+\beta$ ($\alpha\ne0$), so let $M=0$. Then $\E X=\E Z/2$, $\E X^2=\E Z^2/3$, and $X$ has no atom except possibly at $0$.

\emph{Median at the mode.} Suppose $0$ is a median of $X$; this is the situation of (ii) (a symmetric unimodal law has its centre as a mode and a median). Then
$p_\pm=\mathbb P(\pm Z>0)=\mathbb P(\pm X>0)\le\frac12$. With $a=\E Z_+$, $b=\E Z_-$: $\E Z_\pm^2\ge(\E Z_\pm)^2/p_\pm\ge2(\E Z_\pm)^2$ by Cauchy--Schwarz, so
$\mathrm{Var}X\ge\frac23(a^2+b^2)-\frac14(a-b)^2$ and $\tau=\frac{a+b}2>0$; with $a+b=1$, $a=s$, the ratio is at least $4(\frac13s^2-\frac13s+\frac5{12})=\frac43+\frac43(s-\frac12)^2\ge\frac43$.
Equality forces $s=\frac12$ and equality in both Cauchy--Schwarz steps, i.e.\ $\mathbb P(Z>0)=\mathbb P(Z<0)=\frac12$ with $Z$ constant on each event; so $Z=\pm c$ with
probability $\frac12$ each and $X$ is uniform. This proves (ii), and (i) when a mode is a median, since $\frac43>\frac97$.

\emph{Median off the mode.} By reflection and scaling let the median be $1$. Since $X$ has no atom at $1$, $\mathbb P(X\le1)=\frac12$. For a value $z$ of $Z$ put
$\psi(z)=\mathbb P(Uz\le1)$ ($=1$ for $z\le1$, $=1/z$ for $z\ge1$) and $\phi(z)=\E|Uz-1|$ ($=1-z/2$ for $z\le1$, $=z/2-1+1/z$ for $z\ge1$). With $\nu$ the law of $Z$: $\int\psi\,d\nu=\frac12$,
$\tau=\int\phi\,d\nu$, and the claim is
\[H(\nu):=\tfrac73\int z^2d\nu-\tfrac74\Bigl(\int z\,d\nu\Bigr)^2-9\Bigl(\int\phi\,d\nu\Bigr)^2\ge0 .\]
By Lemma~\ref{lem:dec} with $g=\psi$, $\nu=\int\nu_s\,\pi(ds)$ with two-point $\nu_s$ satisfying $\int\psi\,d\nu_s=\frac12$; since $(\int L(\nu_s)\,d\pi)^2\le\int L(\nu_s)^2d\pi$ for linear $L$,
$H(\nu)\ge\int H(\nu_s)\,\pi(ds)$. So let $\nu$ have a left atom $a\le2$ (probability $1-q$) and a right atom $b\ge2$ (probability $q$); the one-point law $Z\equiv2$ gives $H=\frac1{12}$.

\emph{Case $a\le1$.} Here $q=\frac b{2(b-1)}$; $H$ is quadratic in $a$ with leading coefficient $\frac{(b-2)(b+5)}{6(b-1)^2}\ge0$ and minimizer $a^*(b)=-\frac{3(b-3)(b+3)}{4(b+5)}\le1$
(equivalently $(3b+7)(b-1)\ge0$), and
\[\min_aH=\frac{7(b-3)^2(b+6)}{96(b+5)}\ge0,\]
with equality only at $b=3$, $a=0$, $q=\frac34$.

\emph{Case $1\le a\le2$.} Here $q=\frac{b(2-a)}{2(b-a)}$ and $H=N(a,b)/12$ with $N=-12a^2b^2+34a^2b-20a^2+34ab^2-95ab+54a-20b^2+54b-27$. Then $N(a,2)\equiv1$, $N$ is convex
in $b$ (leading coefficient $2(2-a)(6a-5)\ge0$), and $\partial_bN(a,2)=(2-a)(14a-13)\ge0$; hence $N\ge1$.

\emph{Equality} in (i) forces $H(\nu_s)=0$ for $\pi$-a.e.\ $s$, so every component, hence $\nu$, is the law of $Z\in\{0,3\}$ with probabilities $(\frac14,\frac34)$;
then $X$ has an atom $\frac14$ at $0$ and is uniform on $[0,3]$ otherwise, the image of $\frac14\delta_0+\frac34U[0,1]$ under $x\mapsto3x$. Conversely, for
$\frac14\delta_0+\frac34U[0,1]$: median $\frac13$, $\tau=\frac14\cdot\frac13+\frac34\cdot\frac{1/9+4/9}2=\frac7{24}$, $\mathrm{Var}=\frac14-\frac9{64}=\frac7{64}$, ratio $\frac97$.\end{proof}
\begin{corollary}[asymptotically extremal marginal]\label{cor:asym-ext}
For every non-degenerate unimodal $F$ with bounded support, $\lim_{d\to\infty}d\,(2-K_d^F)=\frac{\pi^2}2\frac{\mathrm{Var}X}{\tau^2}\ge\frac{9\pi^2}{14}$, with equality iff $F$ is an
affine image of $\frac14\delta_0+\frac34U[0,1]$; if a median is a mode (in particular for symmetric $F$), the bound is $\frac{2\pi^2}3$, attained by the uniform law.
\end{corollary}
\begin{proof}Cram\'er's condition is automatic: with $p_0=\mathbb P(Z=0)<1$, $\varphi(t)=e^{itM}(p_0+(1-p_0)\psi(t))$, where $\psi$ is the characteristic function of a mixture
of uniform laws, hence absolutely continuous; so $|\varphi(t)|\to p_0<1$, and $|\varphi(t)|<1$ for $t\ne0$ since a law with an absolutely continuous component is not
lattice. Apply Theorem~\ref{thm:price} and Proposition~\ref{prop:97}.\end{proof}
\noindent For finite $d$ we use the fusions $D(t,p)$ of Section~\ref{sec:nonsym} and the mixability criterion (M1) of Section~\ref{sec:prelim}.
\begin{lemma}[atoms at the left end]\label{lem:atom}
Let $G=w_0\delta_0+(1-w_0)H$ with $H$ having a nonincreasing density on $[0,b]$. If the mean of $G$ is at least $b/d$, then $G$ is $d$-completely mixable.
\end{lemma}
\begin{proof}Replace $w_0\delta_0$ by $w_0U[0,\varepsilon]$: the result $G_\varepsilon$ has a nonincreasing density on $[0,b]$ and a larger mean, so it is $d$-CM by (M1). Complete mixes of
$G_\varepsilon$ are tight (support in $[0,b]$); a subsequential limit in law has marginals $G$, and its coordinate sum is the limit of the constants $d\,\E G_\varepsilon$.\end{proof}
\begin{theorem}[finite $d$]\label{thm:finite-d}
Let $F=\frac14\delta_0+\frac34U[0,1]$ and $U$ the uniform law. Then $K_d^F>K_d^U$ for every $d\ge4$. For $d=3$, $F'=\frac{41}{250}\delta_0+\frac{209}{250}U[0,1]$ has
$K_3^{F'}\ge1.0878685$, which exceeds $K_3^U=1.0849633$, $K_3^{F^*}=1.0874849$ (Theorem~\ref{thm:d3}) and the supremum $1.0874848690$ of $\kappa_3^F$ in
Proposition~\ref{prop:d3sup}.
\end{theorem}
\begin{proof}[Proof of Theorem~\ref{thm:finite-d}]\emph{Attainable fusions.} For $F_w=w\delta_0+(1-w)U[0,1]$ and $p>w$, the lower $p$-part is $\frac wp\delta_0+\frac{p-w}pU[0,q]$, $q=\frac{p-w}{1-w}$, with mean $\frac{(1-w)q^2}{2p}$;
the upper part $U[q,1]$ is $d$-CM for every $d\ge2$ by (M1). By Lemma~\ref{lem:atom} the lower part is $d$-CM if $\frac{(1-w)q^2}{2p}\ge\frac qd$. For $w=\frac14$, $p=\frac12$: $q=\frac13$, mean $\frac1{12}\ge\frac1{3d}$ for
$d\ge4$. For $w=\frac{41}{250}$, $p=\frac{123}{250}$, $d=3$: mean $=\frac{82}{627}=\frac q3$ exactly. By Theorem~\ref{thm:fusion-lb}, $K_d^{F_w}\ge D(t,p)$ for all $t>0$.

\emph{$d=3$ and $4\le d\le200$.} $D(t,p)$ is an explicit expression ($\varphi_w(t)=w+(1-w)\frac{e^{it}-1}{it}$, rational $\mu_\pm$) and is bounded below at explicit $t$ by interval
arithmetic; $K_d^U=\kappa_d$ (Corollary~\ref{cor:families}) is bounded above by a grid with a Lipschitz bound (Appendix~\ref{app:unimodal}). All margins are positive (Table~\ref{tab:margins}).

\emph{$d\ge200$.} Let $t_d=\frac{24\pi}{7d}$. Then $dt_d\mu_-=\frac{2\pi}7$ and $dt_d\mu_+=\frac{2\pi}7+2\pi$ ($\mu_-=\frac1{12}$, $\mu_+=\frac23$), so the fusion equals $e^{2\pi i/7}$; and with $\mu=\frac38$,
$\psi(t)=e^{-i\mu t}\varphi(t)$: $\varphi(t_d)^d=e^{9\pi i/7}\psi(t_d)^d=-e^{2\pi i/7}\psi(t_d)^d$. Hence $D(t_d,\frac12)=|1+\psi(t_d)^d|\ge1+\mathrm{Re}\,\psi(t_d)^d$. By Lemma~\ref{lem:tail},
\[D(t_d,\tfrac12)\ge1+e^{-a/d}-\frac{14.6}{d^2},\qquad \kappa_d\le1+e^{-b/d}+\frac{9.98}{d^2},\qquad a=\tfrac{9\pi^2}{14},\ b=\tfrac{2\pi^2}3 .\]
Since $1-e^{-y}\ge ye^{-y}$, $d^2(e^{-a/d}-e^{-b/d})\ge d\,\frac{\pi^2}{42}e^{-b/d}\ge45.4>24.6$ for $d\ge200$.\end{proof}
\begin{table}[ht]\centering\small
\begin{tabular}{rccc}\toprule
$d$ & lower bound for $K_d^F$ & upper bound for $K_d^U$ & margin\\\midrule
3 ($F'$) & 1.0878685 & 1.0849652 & 0.0029 (vs.\ $F^*$: 0.00038)\\
4 & 1.2002921 & 1.1907241 & 0.0096\\
5 & 1.2914829 & 1.2796741 & 0.0118\\
8 & 1.4731083 & 1.4608883 & 0.0122\\
20 & 1.7381118 & 1.7301725 & 0.0079\\
60 & 1.9014265 & 1.8980328 & 0.0034\\
100 & & & 0.0022\\
200 & & & 0.0011\\\bottomrule
\end{tabular}
\caption{Certified margins (interval arithmetic); $F=\frac14\delta_0+\frac34U[0,1]$, median split. The margin behaves like $\frac{\pi^2}{42d}$.}\label{tab:margins}
\end{table}

\section{Open problems}\label{sec:open}
\begin{enumerate}
\item The supremum of $K_3^F$ over symmetric unimodal laws with $3$-CM $|X|$: conjecturally equal to $\sup\kappa_3^F=1.0874848690$ (Proposition~\ref{prop:d3sup}), which
$F^*$ approaches within $1.4\cdot10^{-9}$. Peaked limits of this class (an atom plus a spread-out part, $r=\varphi(t)\le\frac13$) give at most $1+r^3\le28/27$ by Theorem~\ref{thm:profile}.
\item The three-point family $\mathbb P(X=0)=w$, $\mathbb P(X=\pm1)=\frac{1-w}2$ for $d\ge4$: is $K_d=\kappa_d$ exactly when $|X|$ is $d$-CM?
\item The size of the Gaussian gap: numerically $K_3^G\approx1.081<\kappa_3^G=1.0848$.
\item Conjecture~\ref{conj:nonsym}; exact identities $K_d^F=\sup D$ for $d\ge d_F$.
\item Theorem~\ref{thm:price} for non-lattice laws without Cram\'er's condition.
\item The exact value of $K_d^F$ for $F=\frac14\delta_0+\frac34U[0,1]$ (is it $\sup D$?), and the extremal unimodal marginal at a fixed $d$: is it an atom plus a uniform
part for every $d\ge3$?
\end{enumerate}

\appendix
\section{Finite verifications}\label{app:fin}
\subsection{Method}
Each remaining finite claim has the form $F\le\kappa$ on a box, where $F$ is built from elementary functions that are monotone in each variable on the box (or have
explicit Lipschitz bounds). On a sub-box, an upper bound for $F$ is obtained by evaluating every piece at the corner that maximizes it. Sub-boxes are bisected until
each is accepted \emph{by a margin} ($F<\kappa-10^{-3}$) or \emph{by exactness}: in the window, the inequality $\gamma\sin x_0\,\eta^2/x_0\le P(1-\cos\eta)$ on the box, which forces the
second term of $B_d$ to be at most $P_d(x_0)-\cos x_0\le\kappa$; it is checked through $\gamma\sin x_a/(x_aP_d(x_b))\le(1-\cos\eta_b)/\eta_b^2$. Lower bounds for $\kappa$ are the
values of the formula at explicit points. No interval library is used: margins are $\ge10^{-3}$ against rounding errors of order $10^{-15}$. In every box
$\eta\ge\frac\pi2$ is trivial, since then $B_d\le1+P\cos\eta\le1<\kappa$.

\subsection{Two general lemmas}
\begin{lemma}[monotonicity in $d$]\label{lem:mono-d}
If $\log\varphi$ is concave on the first lobe, then $P_d(x)=\varphi(x/(dm))^d$ is nondecreasing in $d\ge3$ for $0\le x\le\pi$ ($\pi/(dm)$ lies in the lobe by Lemma~\ref{lem:lobe}). Consequently
$\kappa_d^F$ is nondecreasing, and exactness on the window, condition (T$_d$) and every inequality of the form $\kappa_d^F\ge c$ persist for all larger $d$ once they hold.
\end{lemma}
\begin{proof}With $\psi=\log\varphi$ concave and $\psi(0)=0$, $\psi(u)/u$ is nonincreasing, so $d\,\psi(x/(dm))=\frac xm\cdot\frac{\psi(x/(dm))}{x/(dm)}$ is nondecreasing in $d$. Exactness
is monotone in $P$.\end{proof}
\begin{lemma}[the region $x_0<x_c$ for large $d$]\label{lem:binf}
If $\gamma\le\gamma_0$, then $B_d(x_0,\eta)\le1-\cos x_c=x_c\sin x_c=1.68915\ldots$ for all $0<x_0\le x_c$ and $\eta\in[0,\pi]$, whatever $P_d(x_0)\in[0,1]$. Hence this part of (W$_d$)
holds as soon as $\kappa_d^F\ge1-\cos x_c$.
\end{lemma}
\begin{proof}For $x_0\le x_c$ the second term is $-1+\sin x_c\,(x_0+e)+P\cos\eta$ with $e=\min\{\gamma\eta^2/x_0,\eta\}$, and $x_0+e\le\max\{x_c+\gamma\eta^2/x_c,\ (1+\gamma)\eta\}$. If the first
entry is the larger, the second term is at most $x_c\sin x_c+\gamma\sin x_c\,\eta^2/x_c-(1-\cos\eta)\le x_c\sin x_c$, since $(1-\cos\eta)/\eta^2\ge2/\pi^2$ and $\gamma\le\gamma_0$ (for $\cos\eta<0$
the first term is already $\le1$). Otherwise $\eta\ge x_c/(1+\gamma)\ge1.41$ and the first term is at most $1+\cos1.41<1.16$.\end{proof}

\subsection{The verifications, law by law}
\begin{center}\footnotesize\setlength{\tabcolsep}{3pt}
\begin{tabular}{p{2.1cm}p{2.3cm}p{2.9cm}p{3.3cm}p{2.6cm}}\toprule
law & $\log\varphi$ concave & tail & window $[x_c,\pi]$ & $x_0<x_c$\\\midrule
uniform & analytic ($\sinc$) & (T$_d$) for $d\ge5$; $d=4$: Lemma~\ref{lem:onebeyond} ($\beta_{\max}=0.762<\beta_1=1.243$); $d=3$: Lemma~\ref{lem:c3} & exactness at $d=3$ (C1), then all $d$ by Lemma~\ref{lem:mono-d} & boxes $3\le d\le16$ (5--25 each, margin $\ge0.005$); $d\ge17$ by Lemma~\ref{lem:binf}\\
triangular & analytic ($\sinc^2$) & (T$_d$): $1+0.0472<1.0766$ & boxes $3\le d\le19$; exactness at $d=20$ (ratio $0.273$), then Lemma~\ref{lem:mono-d} & boxes $3\le d\le18$; $d\ge19$ by Lemma~\ref{lem:binf}\\
cubic B-spline & analytic ($\sinc^3$) & (T$_d$): $s^*\le0.0103$ & boxes $d=3,4,5$; exactness for $d\ge6$ & boxes $3\le d\le19$; $d\ge20$ by Lemma~\ref{lem:binf}\\
Epanechnikov & brackets ($\varphi\varphi''-\varphi'^2\le-0.021$) & (T$_d$) for $d\ge4$: $s^*\le0.0862$ & boxes $d=4$; exactness for $d\ge5$ & boxes $4\le d\le17$; $d\ge18$ by Lemma~\ref{lem:binf}\\
$F^*$ ($d=3$) & brackets ($\le-0.949$ on $[0,0.599)$) & Lemma~\ref{lem:onebeyond}: $\beta_{\max}=1.0891<\beta_1=1.0932$; $s^{*2}<\tau$ & 25 boxes (18 by exactness) & 29 boxes\\
Gaussian & analytic & none ($s^*=0$) & boxes $d=3,4,5$; exactness for $d\ge6$ (ratio $0.889$ at $d=6$); strict at $x_0^*$ ($0.168,0.126,0.101<4/\pi^2$) & boxes $3\le d\le19$; $d\ge20$ by Lemma~\ref{lem:binf}\\\bottomrule
\end{tabular}\end{center}
\noindent ``Exactness for $d\ge d_1$'' uses the crude form $\gamma\frac{\sin x_c}{x_c}\le\frac4{\pi^2}P_d(\pi)$, i.e.\ exactness on the whole window for $\eta\le\frac\pi2$ at once, then Lemma~\ref{lem:mono-d}.
Mixability for the lower bound: uniform and triangular $d\ge3$ ($b/m=2,3$), B-spline $d\ge4$ ($b/m=48/13$), Epanechnikov $d\ge3$ ($b/m=8/3$), $F^*$ ($b/m=3$).

\subsection{The uniform tail for $d=3$}
\begin{lemma}\label{lem:c3}
For uniform marginals on $[-1,1]$ and $d=3$, if exactly one frequency $t_1\in[\pi,3\pi]$ lies beyond the lobe and $\beta=\frac12(t_2+t_3)\in[0.993,1.599]$, the defect is at most $\kappa_3$.
\end{lemma}
\begin{proof}[Proof sketch]In every direction the defect is at most $1-\frac2{\pi^2}(\rho(\psi)-\beta)_+^2+Q\cos u$, where $\rho(\psi)=\E\,\mathrm{dist}(t_1X_1-\psi,2\pi\mathbb Z)$ (Lemma~\ref{lem:latsep} and
Jensen), $u$ is the distance of $\psi$ from the dangerous direction and $Q\ge|P|$. Only $u\le\arccos(\tau/Q)$ matters; $\partial_\psi\rho=0$ at the dangerous direction by symmetry and
$|\partial_\psi^2\rho|\le1/t_1\le1/\pi$, so $\rho\ge\rho^*-u^2/(2\pi)$ with $\rho^*$ in closed form; the resulting lower bound is linear in $v=u^2$, so two closed-form inequalities at
$v=0$ and $v=u_0^2$ suffice. On 40 cells in $(t_1,\beta)$, with bounds from the monotonicity of $\rho^*$ (quasi-concave), $|\sinc|$ (unimodal on each lobe) and $\sinc(\cdot)^2$,
both inequalities hold with slack $\ge0.0014$ (table in the supplementary material). Below $\beta=0.993$ Lemma~\ref{lem:onebeyond} applies; above $1.599$, $|P|\le\tau$. For $t_1\ge3\pi$, $|\varphi(t_1)|\le1/(3\pi)$, so $|P|>\tau$ forces $\beta<0.81$, while
Lemma~\ref{lem:onebeyond} settles $\beta\le1.12$.\end{proof}

\subsection{Reproducibility}
All certificates run in about $8$ seconds (\texttt{run\_all\_certificates.py}, all passing); the $40$-cell table and the scripts for $F^*$ and the Gaussian law are in the supplementary material.

\subsection{Certification for Section~\ref{sec:unimodal}}\label{app:unimodal}
Unlike the rest of this appendix, the comparison of Theorem~\ref{thm:finite-d} uses interval arithmetic. \emph{Lower bounds.} $D(t,p)$ is evaluated with real interval arithmetic (mpmath, 40 digits): $\varphi_w(t)=w+(1-w)\frac{\sin t}t+i(1-w)\frac{1-\cos t}t$, the power $\varphi^d$ by repeated
complex multiplication of intervals, and $|\cdot|$ by the square root of a sum of squares; the lower endpoint is reported.

\emph{Upper bounds for $\kappa_d=\max_x\{\sinc(2x/d)^d-\cos x\}$.} On $[0,\pi]$: a grid of step $h=2\cdot10^{-6}$ ($d\le60$) or $10^{-6}$ ($61\le d\le200$) and $|g'|\le2\max|\sinc'|+1<1.88$. For $x>\pi$: $|\sinc(2x/d)|\le\sinc(2\pi/d)$ while
$2x/d\le\pi$, and $|\sinc|\le0.2173$ beyond, so the value is at most $1+\sinc(2\pi/d)^d$, the value at $\pi$.
\begin{lemma}[tail constants, $d\ge200$]\label{lem:tail}
With $t_d=\frac{24\pi}{7d}$: $D(t_d,\frac12)\ge1+e^{-a/d}-14.6/d^2$ and $\kappa_d\le1+e^{-b/d}+9.98/d^2$.
\end{lemma}
\begin{proof}\emph{Lower.} $Y=X-\frac38\in[-\frac38,\frac58]$, $\E Y^2=\sigma^2=\frac7{64}$, $\E|Y|^3\le\frac58\sigma^2$. From $|e^{iy}-1-iy+\frac{y^2}2|\le\frac{|y|^3}6$: $\psi(t)=1-\frac{\sigma^2t^2}2+R$,
$|R|\le r:=\frac{35}{3072}|t|^3$. With $u=\psi-1$, $|u|\le\frac{\sigma^2t^2}2+r\le\frac12$ and $|\log(1+u)-u|\le|u|^2$, so $d\log\psi(t_d)=-\frac ad+de_1$ with $d|e_1|\le d(r+|u|^2)\le\frac{14.24}{d^2}+\frac{41.2}{d^3}\le\frac{14.45}{d^2}$
($\frac{\sigma^2t_d^2}2=\frac a{d^2}$, $dr=14.24/d^2$). Then $\mathrm{Re}\,\psi^d\ge e^{-a/d-\epsilon}\cos\epsilon\ge e^{-a/d}(1-\epsilon-\epsilon^2/2)$ with $\epsilon=d|e_1|$.

\emph{Upper.} On $[0,\pi]$, $\sinc y\le e^{-y^2/6}$ ($|y|<\pi$; $\log\sinc y=\sum_k\log(1-\frac{y^2}{k^2\pi^2})\le-\frac{y^2}6$) gives $g(x)\le e^{-\alpha x^2}-\cos x$, $\alpha=\frac2{3d}$. For $x=\pi-\delta$,
$\delta\le1$: $e^{-\alpha(\pi-\delta)^2}\le e^{-\alpha\pi^2}+2\alpha\pi e^{2\alpha\pi}\delta$ and $-\cos x=\cos\delta\le1-\frac{11}{24}\delta^2$, so $g\le1+e^{-b/d}+\frac6{11}(2\alpha\pi e^{2\alpha\pi})^2\le1+e^{-b/d}+9.98/d^2$. For
$\delta\ge1$, $g\le1+\cos1<1.55<1+e^{-b/d}$; for $x>\pi$ see above.\end{proof}

\section*{Use of AI tools}
AI systems used in the research and manuscript workflow included OpenAI GPT-5.6 Sol and Anthropic Claude Opus 5.5. They were used for
mathematical exploration, proof development and critique, symbolic and numerical verification, literature search, and manuscript drafting
and editing. AI systems are not authors. The author reviewed the mathematical statements and proofs and takes full responsibility for the
content.

\begin{thebibliography}{9}
\bibitem{I} A.~Martinevski, A sharp universal profile for multivariate characteristic-function factorization, preprint (2026).
\bibitem{WW11} B.~Wang, R.~Wang, The complete mixability and convex minimization problems with monotone marginal densities,
\emph{J.\ Multivariate Anal.} 102 (2011) 1344--1360.
\bibitem{PWW12} G.~Puccetti, B.~Wang, R.~Wang, Advances in complete mixability, \emph{J.\ Appl.\ Probab.} 49 (2012) 430--440.
\bibitem{MN79} I.~Meilijson, A.~N\'adas, Convex majorization with an application to the length of critical paths, \emph{J.\ Appl.\ Probab.} 16 (1979) 671--677.
\bibitem{Dh02} J.~Dhaene, M.~Denuit, M.~J.~Goovaerts, R.~Kaas, D.~Vyncke, The concept of comonotonicity in actuarial science and finance: theory,
\emph{Insurance Math.\ Econom.} 31 (2002) 3--33, doi:10.1016/S0167-6687(02)00134-8.
\bibitem{WW16} B.~Wang, R.~Wang, Joint mixability, \emph{Math.\ Oper.\ Res.} 41 (2016) 808--826.
\bibitem{DJ88} S.~Dharmadhikari, K.~Joag-Dev, \emph{Unimodality, Convexity, and Applications}, Academic Press, 1988.
\bibitem{Kh38} A.~Ya.~Khinchine, On unimodal distributions, \emph{Izv.\ Nauchno-Issled.\ Inst.\ Mat.\ Mekh.\ Tomsk.\ Gos.\ Univ.} 2 (1938) 1--7.
\bibitem{LSWY18} L.~Li, H.~Shao, R.~Wang, J.~Yang, Worst-case Range Value-at-Risk with partial information, \emph{SIAM J.\ Financial Math.} 9 (2018) 190--218.
\bibitem{BKV20} C.~Bernard, R.~Kazzi, S.~Vanduffel, Range Value-at-Risk bounds for unimodal distributions under partial information, \emph{Insurance Math.\ Econom.} 94 (2020)
9--24; Corrigendum and addendum, \emph{ibid.} 112 (2023) 110--119.
\bibitem{ZBYS26} M.~Zhao, N.~Balakrishnan, C.~Yin, H.~Shao, Extremal cases of distortion risk measures with partial information, \emph{Stat.\ Methods Appl.} (2026),
doi:10.1007/s10260-026-00867-8.
\bibitem{ZY25} B.~Zuo, C.~Yin, Analyzing distortion riskmetrics and weighted entropy for unimodal and symmetric distributions under partial information constraints,
arXiv:2504.19725 (2025).
\bibitem{Wi1866} A.~Winckler, Allgemeine S\"atze zur Theorie der unregelm\"assigen Beobachtungsfehler, \emph{Sitzungsber.\ Math.-Naturwiss.\ Kl.\ Kais.\ Akad.\ Wiss.\ Wien}
53 (1866) 6--41.
\end{thebibliography}
\end{document}
